\subsection*{Immersion / Comprehensible Input}
\textbf{Qué es:} Exponerte a contenido que entiendes en un 70--90\%, manteniendo flujo natural.

\textbf{Para qué sirve:} El cerebro detecta patrones y estructura sin sobrecarga; reduce la traducción palabra por palabra.

\textbf{Cómo aplicarla:}
\begin{itemize}[leftmargin=*]
	\item Empieza con videos/podcasts A1 o canales lentos; activa subtítulos en neerlandés si ayudan.
	\item Haz sesiones cortas (10--15 min) diarias: una escucha para idea general, otra con pausa ligera para notar frases clave.
	\item Marca 3--5 expresiones útiles por sesión; no subrayes todo.
\end{itemize}

\textbf{Ejemplos de materiales:}
\begin{itemize}[leftmargin=*]
	\item Noticias lentas para aprendices (\textit{Nieuwsbegrip}, podcasts A1).
	\item Videos de recetas simples o vlogs con lenguaje cotidiano.
\end{itemize}
